\documentclass[10pt]{article}
% Hojas resumen: un A4 por delante y por detrás, dos columnas por cara.
% Uso: \documentclass[10pt]{article} % Hojas resumen: un A4 por delante y por detrás, dos columnas por cara.
% Uso: \documentclass[10pt]{article} % Hojas resumen: un A4 por delante y por detrás, dos columnas por cara.
% Uso: \documentclass[10pt]{article} \input{estilo-resumen.tex} \bloque{analisis|algebra|azar} \cabecera{n}{Título}{Lecciones}
\usepackage[utf8]{inputenc}
\usepackage[T1]{fontenc}
\usepackage[spanish,es-nodecimaldot]{babel}
\usepackage{lmodern}
\usepackage{amsmath,amssymb}
\usepackage{graphicx}
\usepackage[a4paper,left=10mm,right=10mm,top=15mm,bottom=13mm,headheight=11pt,headsep=3mm,footskip=6mm]{geometry}
\usepackage{multicol}
\setlength{\columnsep}{6mm}
\setlength{\columnseprule}{0.3pt}
\usepackage{xcolor}
\usepackage{enumitem}
\setlist{nosep,leftmargin=1.1em,itemsep=1pt}
\usepackage{array}
\usepackage{needspace}
\usepackage{pgfplots}
\usepgfplotslibrary{fillbetween}
\pgfplotsset{compat=1.18,
  mini/.style={width=4.3cm,height=3.3cm,scale only axis=false,axis lines=middle,xtick=\empty,ytick=\empty,samples=120,
    xlabel={\scriptsize$x$},ylabel={\scriptsize$y$},title style={font=\scriptsize,yshift=-4pt},clip=true,
    every axis plot/.append style={thick,color=acento}}}
\usepackage{ragged2e}
\usepackage[most]{tcolorbox}
\usepackage{fancyhdr,lastpage}
\usepackage[hidelinks,pdfencoding=auto]{hyperref}
\hypersetup{pdflang=es-ES,pdfsubject={Matemáticas Aplicadas a las Ciencias Sociales II}}

\definecolor{analisis}{HTML}{1D4ED8}
\definecolor{algebra}{HTML}{86198F}
\definecolor{azar}{HTML}{9A3412}
\definecolor{aviso}{HTML}{B91C1C}
\colorlet{acento}{analisis}
\newcommand{\bloque}[1]{\colorlet{acento}{#1}}

\setlength{\parindent}{0pt}
\setlength{\parskip}{2pt}
\renewcommand{\arraystretch}{1.2}
\AtBeginDocument{\RaggedRight\fontsize{9.8}{11.9}\selectfont\setlength{\abovedisplayskip}{3pt}\setlength{\belowdisplayskip}{3pt}\setlength{\abovedisplayshortskip}{1pt}\setlength{\belowdisplayshortskip}{1pt}}

\newcommand{\temanum}{}\newcommand{\tematitulo}{}
\pagestyle{fancy}
\fancyhf{}
\fancyhead[L]{\footnotesize\color{acento}\textbf{Hoja resumen}\quad Tema \temanum. \tematitulo}
\fancyhead[R]{\footnotesize Matemáticas Aplicadas a las Ciencias Sociales II --- 2.\textsuperscript{o} Bachillerato}
\fancyfoot[L]{\scriptsize IES Al-Ándalus --- Juan María Gámez Ortiz, con la ayuda de Claude y de Codex --- CC BY-NC 4.0}
\fancyfoot[R]{\scriptsize Cara \thepage\ de \pageref{LastPage}}
\renewcommand{\headrulewidth}{0.4pt}

% Cabecera de la primera cara: número, título y lecciones del tema.
\newcommand{\cabecera}[3]{%
  \renewcommand{\temanum}{#1}\renewcommand{\tematitulo}{#2}%
  \hypersetup{pdftitle={Hoja resumen. Tema #1. #2}}%
  \noindent\begin{tcolorbox}[enhanced,colback=acento!8,colframe=acento,boxrule=0.6pt,arc=2pt,left=4pt,right=4pt,top=2pt,bottom=2pt]
  \begin{minipage}[c]{1.1cm}\includegraphics[width=1cm]{logo-al-andalus.png}\end{minipage}%
  \begin{minipage}[c]{\dimexpr\linewidth-1.2cm}{\large\bfseries\color{acento}Tema #1. #2}\\[1pt]{\scriptsize #3}\end{minipage}
  \end{tcolorbox}\vspace{-2pt}}

% Apartado con filete del color del bloque.
\newcommand{\seccion}[1]{\par\needspace{5\baselineskip}\vspace{3pt}{\color{acento}\bfseries\normalsize #1}\par\vspace{-1pt}{\color{acento}\hrule height 0.5pt}\vspace{2pt}}
\newcommand{\sub}[1]{\par\needspace{3\baselineskip}\vspace{1pt}\textbf{\color{acento}#1}\enspace}

% Fórmula clave, aviso de errores y ejemplo.
\tcbset{before upper=\RaggedRight}
\newtcolorbox{formula}{enhanced,colback=acento!6,colframe=acento!70,boxrule=0.5pt,arc=1.5pt,left=3pt,right=3pt,top=0pt,bottom=0pt,before skip=2pt,after skip=2pt}
\newtcolorbox{ojo}{enhanced,colback=aviso!5,colframe=aviso!60,boxrule=0.5pt,arc=1.5pt,left=3pt,right=3pt,top=1pt,bottom=1pt,before skip=3pt,after skip=2pt,
  fonttitle=\bfseries\footnotesize,coltitle=aviso,colbacktitle=aviso!5,title={Ojo},attach title to upper={\enspace}}
\newtcolorbox{ejemplo}{enhanced,colback=white,colframe=acento!45,boxrule=0.4pt,arc=1.5pt,left=3pt,right=3pt,top=1pt,bottom=1pt,before skip=3pt,after skip=2pt,
  fonttitle=\bfseries\footnotesize,coltitle=acento,colbacktitle=white,title={Ejemplo},attach title to upper={\enspace}}
\newcommand{\dx}{\,dx}
 \bloque{analisis|algebra|azar} \cabecera{n}{Título}{Lecciones}
\usepackage[utf8]{inputenc}
\usepackage[T1]{fontenc}
\usepackage[spanish,es-nodecimaldot]{babel}
\usepackage{lmodern}
\usepackage{amsmath,amssymb}
\usepackage{graphicx}
\usepackage[a4paper,left=10mm,right=10mm,top=15mm,bottom=13mm,headheight=11pt,headsep=3mm,footskip=6mm]{geometry}
\usepackage{multicol}
\setlength{\columnsep}{6mm}
\setlength{\columnseprule}{0.3pt}
\usepackage{xcolor}
\usepackage{enumitem}
\setlist{nosep,leftmargin=1.1em,itemsep=1pt}
\usepackage{array}
\usepackage{needspace}
\usepackage{pgfplots}
\usepgfplotslibrary{fillbetween}
\pgfplotsset{compat=1.18,
  mini/.style={width=4.3cm,height=3.3cm,scale only axis=false,axis lines=middle,xtick=\empty,ytick=\empty,samples=120,
    xlabel={\scriptsize$x$},ylabel={\scriptsize$y$},title style={font=\scriptsize,yshift=-4pt},clip=true,
    every axis plot/.append style={thick,color=acento}}}
\usepackage{ragged2e}
\usepackage[most]{tcolorbox}
\usepackage{fancyhdr,lastpage}
\usepackage[hidelinks,pdfencoding=auto]{hyperref}
\hypersetup{pdflang=es-ES,pdfsubject={Matemáticas Aplicadas a las Ciencias Sociales II}}

\definecolor{analisis}{HTML}{1D4ED8}
\definecolor{algebra}{HTML}{86198F}
\definecolor{azar}{HTML}{9A3412}
\definecolor{aviso}{HTML}{B91C1C}
\colorlet{acento}{analisis}
\newcommand{\bloque}[1]{\colorlet{acento}{#1}}

\setlength{\parindent}{0pt}
\setlength{\parskip}{2pt}
\renewcommand{\arraystretch}{1.2}
\AtBeginDocument{\RaggedRight\fontsize{9.8}{11.9}\selectfont\setlength{\abovedisplayskip}{3pt}\setlength{\belowdisplayskip}{3pt}\setlength{\abovedisplayshortskip}{1pt}\setlength{\belowdisplayshortskip}{1pt}}

\newcommand{\temanum}{}\newcommand{\tematitulo}{}
\pagestyle{fancy}
\fancyhf{}
\fancyhead[L]{\footnotesize\color{acento}\textbf{Hoja resumen}\quad Tema \temanum. \tematitulo}
\fancyhead[R]{\footnotesize Matemáticas Aplicadas a las Ciencias Sociales II --- 2.\textsuperscript{o} Bachillerato}
\fancyfoot[L]{\scriptsize IES Al-Ándalus --- Juan María Gámez Ortiz, con la ayuda de Claude y de Codex --- CC BY-NC 4.0}
\fancyfoot[R]{\scriptsize Cara \thepage\ de \pageref{LastPage}}
\renewcommand{\headrulewidth}{0.4pt}

% Cabecera de la primera cara: número, título y lecciones del tema.
\newcommand{\cabecera}[3]{%
  \renewcommand{\temanum}{#1}\renewcommand{\tematitulo}{#2}%
  \hypersetup{pdftitle={Hoja resumen. Tema #1. #2}}%
  \noindent\begin{tcolorbox}[enhanced,colback=acento!8,colframe=acento,boxrule=0.6pt,arc=2pt,left=4pt,right=4pt,top=2pt,bottom=2pt]
  \begin{minipage}[c]{1.1cm}\includegraphics[width=1cm]{logo-al-andalus.png}\end{minipage}%
  \begin{minipage}[c]{\dimexpr\linewidth-1.2cm}{\large\bfseries\color{acento}Tema #1. #2}\\[1pt]{\scriptsize #3}\end{minipage}
  \end{tcolorbox}\vspace{-2pt}}

% Apartado con filete del color del bloque.
\newcommand{\seccion}[1]{\par\needspace{5\baselineskip}\vspace{3pt}{\color{acento}\bfseries\normalsize #1}\par\vspace{-1pt}{\color{acento}\hrule height 0.5pt}\vspace{2pt}}
\newcommand{\sub}[1]{\par\needspace{3\baselineskip}\vspace{1pt}\textbf{\color{acento}#1}\enspace}

% Fórmula clave, aviso de errores y ejemplo.
\tcbset{before upper=\RaggedRight}
\newtcolorbox{formula}{enhanced,colback=acento!6,colframe=acento!70,boxrule=0.5pt,arc=1.5pt,left=3pt,right=3pt,top=0pt,bottom=0pt,before skip=2pt,after skip=2pt}
\newtcolorbox{ojo}{enhanced,colback=aviso!5,colframe=aviso!60,boxrule=0.5pt,arc=1.5pt,left=3pt,right=3pt,top=1pt,bottom=1pt,before skip=3pt,after skip=2pt,
  fonttitle=\bfseries\footnotesize,coltitle=aviso,colbacktitle=aviso!5,title={Ojo},attach title to upper={\enspace}}
\newtcolorbox{ejemplo}{enhanced,colback=white,colframe=acento!45,boxrule=0.4pt,arc=1.5pt,left=3pt,right=3pt,top=1pt,bottom=1pt,before skip=3pt,after skip=2pt,
  fonttitle=\bfseries\footnotesize,coltitle=acento,colbacktitle=white,title={Ejemplo},attach title to upper={\enspace}}
\newcommand{\dx}{\,dx}
 \bloque{analisis|algebra|azar} \cabecera{n}{Título}{Lecciones}
\usepackage[utf8]{inputenc}
\usepackage[T1]{fontenc}
\usepackage[spanish,es-nodecimaldot]{babel}
\usepackage{lmodern}
\usepackage{amsmath,amssymb}
\usepackage{graphicx}
\usepackage[a4paper,left=10mm,right=10mm,top=15mm,bottom=13mm,headheight=11pt,headsep=3mm,footskip=6mm]{geometry}
\usepackage{multicol}
\setlength{\columnsep}{6mm}
\setlength{\columnseprule}{0.3pt}
\usepackage{xcolor}
\usepackage{enumitem}
\setlist{nosep,leftmargin=1.1em,itemsep=1pt}
\usepackage{array}
\usepackage{needspace}
\usepackage{pgfplots}
\usepgfplotslibrary{fillbetween}
\pgfplotsset{compat=1.18,
  mini/.style={width=4.3cm,height=3.3cm,scale only axis=false,axis lines=middle,xtick=\empty,ytick=\empty,samples=120,
    xlabel={\scriptsize$x$},ylabel={\scriptsize$y$},title style={font=\scriptsize,yshift=-4pt},clip=true,
    every axis plot/.append style={thick,color=acento}}}
\usepackage{ragged2e}
\usepackage[most]{tcolorbox}
\usepackage{fancyhdr,lastpage}
\usepackage[hidelinks,pdfencoding=auto]{hyperref}
\hypersetup{pdflang=es-ES,pdfsubject={Matemáticas Aplicadas a las Ciencias Sociales II}}

\definecolor{analisis}{HTML}{1D4ED8}
\definecolor{algebra}{HTML}{86198F}
\definecolor{azar}{HTML}{9A3412}
\definecolor{aviso}{HTML}{B91C1C}
\colorlet{acento}{analisis}
\newcommand{\bloque}[1]{\colorlet{acento}{#1}}

\setlength{\parindent}{0pt}
\setlength{\parskip}{2pt}
\renewcommand{\arraystretch}{1.2}
\AtBeginDocument{\RaggedRight\fontsize{9.8}{11.9}\selectfont\setlength{\abovedisplayskip}{3pt}\setlength{\belowdisplayskip}{3pt}\setlength{\abovedisplayshortskip}{1pt}\setlength{\belowdisplayshortskip}{1pt}}

\newcommand{\temanum}{}\newcommand{\tematitulo}{}
\pagestyle{fancy}
\fancyhf{}
\fancyhead[L]{\footnotesize\color{acento}\textbf{Hoja resumen}\quad Tema \temanum. \tematitulo}
\fancyhead[R]{\footnotesize Matemáticas Aplicadas a las Ciencias Sociales II --- 2.\textsuperscript{o} Bachillerato}
\fancyfoot[L]{\scriptsize IES Al-Ándalus --- Juan María Gámez Ortiz, con la ayuda de Claude y de Codex --- CC BY-NC 4.0}
\fancyfoot[R]{\scriptsize Cara \thepage\ de \pageref{LastPage}}
\renewcommand{\headrulewidth}{0.4pt}

% Cabecera de la primera cara: número, título y lecciones del tema.
\newcommand{\cabecera}[3]{%
  \renewcommand{\temanum}{#1}\renewcommand{\tematitulo}{#2}%
  \hypersetup{pdftitle={Hoja resumen. Tema #1. #2}}%
  \noindent\begin{tcolorbox}[enhanced,colback=acento!8,colframe=acento,boxrule=0.6pt,arc=2pt,left=4pt,right=4pt,top=2pt,bottom=2pt]
  \begin{minipage}[c]{1.1cm}\includegraphics[width=1cm]{logo-al-andalus.png}\end{minipage}%
  \begin{minipage}[c]{\dimexpr\linewidth-1.2cm}{\large\bfseries\color{acento}Tema #1. #2}\\[1pt]{\scriptsize #3}\end{minipage}
  \end{tcolorbox}\vspace{-2pt}}

% Apartado con filete del color del bloque.
\newcommand{\seccion}[1]{\par\needspace{5\baselineskip}\vspace{3pt}{\color{acento}\bfseries\normalsize #1}\par\vspace{-1pt}{\color{acento}\hrule height 0.5pt}\vspace{2pt}}
\newcommand{\sub}[1]{\par\needspace{3\baselineskip}\vspace{1pt}\textbf{\color{acento}#1}\enspace}

% Fórmula clave, aviso de errores y ejemplo.
\tcbset{before upper=\RaggedRight}
\newtcolorbox{formula}{enhanced,colback=acento!6,colframe=acento!70,boxrule=0.5pt,arc=1.5pt,left=3pt,right=3pt,top=0pt,bottom=0pt,before skip=2pt,after skip=2pt}
\newtcolorbox{ojo}{enhanced,colback=aviso!5,colframe=aviso!60,boxrule=0.5pt,arc=1.5pt,left=3pt,right=3pt,top=1pt,bottom=1pt,before skip=3pt,after skip=2pt,
  fonttitle=\bfseries\footnotesize,coltitle=aviso,colbacktitle=aviso!5,title={Ojo},attach title to upper={\enspace}}
\newtcolorbox{ejemplo}{enhanced,colback=white,colframe=acento!45,boxrule=0.4pt,arc=1.5pt,left=3pt,right=3pt,top=1pt,bottom=1pt,before skip=3pt,after skip=2pt,
  fonttitle=\bfseries\footnotesize,coltitle=acento,colbacktitle=white,title={Ejemplo},attach title to upper={\enspace}}
\newcommand{\dx}{\,dx}

\bloque{analisis}
\begin{document}
\cabecera{3}{Integrales}{Lecciones 3.1--3.8: integral indefinida, inmediatas, casi inmediatas, métodos de integración, integral definida, teorema fundamental del cálculo, regla de Barrow y aplicaciones al cálculo de áreas.}
\begin{multicols}{2}

\seccion{1. Primitiva e integral indefinida}
\(F\) es una \textbf{primitiva} de \(f\) si \(F'(x)=f(x)\). Si \(F\) es una, todas son \(F+C\).
\begin{formula}
\[\int f(x)\,dx=F(x)+C,\qquad C\in\mathbb R\]
\end{formula}
\textbf{Comprobar siempre derivando}: la derivada del resultado debe ser el integrando.
\sub{Linealidad.}
\[\int\bigl(\alpha f+\beta g\bigr)\,dx=\alpha\!\int f\,dx+\beta\!\int g\,dx\]
No hay regla para productos ni cocientes: \(\int fg\,dx\ne\int f\,dx\cdot\int g\,dx\).
\begin{ejemplo}
Primitiva de \(f(x)=6x-4\) con \(F(1)=3\): \(F(x)=3x^2-4x+C\), \(F(1)=-1+C=3\Rightarrow C=4\). \(F(x)=3x^2-4x+4\).
\end{ejemplo}

\seccion{2. Tabla de integrales}
\(u=u(x)\) es la función interior y \(u'\) su derivada; la forma compuesta exige que \(u'\) aparezca como factor (salvo una constante).
\begin{center}\renewcommand{\arraystretch}{1.75}\small
\begin{tabular}{@{}l@{\ \ }l@{}}
\textbf{Inmediata} & \textbf{Compuesta}\\ \hline
\(\displaystyle\int k\,dx=kx+C\) & \\
\(\displaystyle\int x^n\,dx=\frac{x^{n+1}}{n+1}+C\) & \(\displaystyle\int u'u^n\,dx=\frac{u^{n+1}}{n+1}+C\)\\
\multicolumn{2}{@{}l}{\scriptsize (\(n\ne-1\); incluye raíces: \(\sqrt x=x^{1/2}\), \(\frac1{\sqrt x}=x^{-1/2}\))}\\
\(\displaystyle\int\frac1x\,dx=\ln|x|+C\) & \(\displaystyle\int\frac{u'}u\,dx=\ln|u|+C\)\\
\(\displaystyle\int e^x\,dx=e^x+C\) & \(\displaystyle\int u'e^u\,dx=e^u+C\)\\
\(\displaystyle\int a^x\,dx=\frac{a^x}{\ln a}+C\) & \(\displaystyle\int u'a^u\,dx=\frac{a^u}{\ln a}+C\)\\
\(\displaystyle\int\cos x\,dx=\sin x+C\) & \(\displaystyle\int u'\cos u\,dx=\sin u+C\)\\
\(\displaystyle\int\sin x\,dx=-\cos x+C\) & \(\displaystyle\int u'\sin u\,dx=-\cos u+C\)
\end{tabular}
\end{center}
\begin{formula}
\textbf{Caso lineal} \(u=ax+b\), \(a\ne0\): se divide entre \(a\).
\[\int(ax+b)^n\,dx=\frac{(ax+b)^{n+1}}{a(n+1)}+C\]
\[\int\frac{dx}{ax+b}=\frac1a\ln|ax+b|+C,\qquad \int e^{ax+b}\,dx=\frac{e^{ax+b}}{a}+C\]
\end{formula}

\seccion{3. Integrales casi inmediatas}
\begin{enumerate}
  \item Identifica la función interior \(u\) y calcula \(u'\).
  \item Si falta un \textbf{factor numérico}, multiplica y divide por él.
  \item Aplica la fila de la tabla y comprueba derivando.
\end{enumerate}
\begin{ejemplo}
\[\int x(x^2+1)^3\,dx=\frac12\int2x(x^2+1)^3\,dx=\frac{(x^2+1)^4}{8}+C\]
\[\int\frac{5x}{x^2+4}\,dx=\frac52\ln(x^2+4)+C\]
\[\int e^{2-3x}\,dx=-\frac13e^{2-3x}+C\]
\end{ejemplo}
\begin{ojo}
Solo se ajustan \textbf{números}. \(\int e^{x^2}\,dx\) no es casi inmediata: falta el factor \(x\), que no es constante.
\end{ojo}

\seccion{4. Métodos de integración}
\sub{Simplificar antes.} Desarrollar productos, dividir cada término entre el denominador, factorizar y cancelar (conservando el dominio).
\begin{ejemplo}
\[\int\frac{x^2+1}{x+1}\,dx=\int\Bigl(x-1+\frac2{x+1}\Bigr)dx\]
\[=\frac{x^2}2-x+2\ln|x+1|+C\]
\end{ejemplo}
\sub{Cambio de variable.}
\begin{formula}
\[\int h(g(x))\,g'(x)\,dx=\int h(u)\,du=H(u)+C=H(g(x))+C\]
\end{formula}
Se escribe toda la integral en \(u\), incluido \(du\), y al final se vuelve a \(x\).
\begin{ejemplo}
\(\int x\sqrt{x^2+4}\,dx\): \(u=x^2+4\), \(du=2x\,dx\).
\[\frac12\int u^{1/2}\,du=\frac12\cdot\frac{u^{3/2}}{3/2}+C=\frac13(x^2+4)^{3/2}+C\]
\end{ejemplo}
\sub{Por partes} (ampliación): \(\int u\,dv=uv-\int v\,du\). Por ejemplo, \(\int xe^x\,dx=(x-1)e^x+C\).

\seccion{5. Integral definida}
\(\int_a^bf(x)\,dx\) es el límite de las sumas de áreas de rectángulos (sumas de Riemann). Si \(f\ge0\) en \([a,b]\), es el área bajo la curva; si \(f\le0\), es el área con signo negativo.
\sub{Propiedades.}
\[\int_a^af=0,\qquad \int_b^af=-\int_a^bf,\qquad \int_a^bf=\int_a^cf+\int_c^bf\]
\[\int_a^b(\alpha f+\beta g)=\alpha\int_a^bf+\beta\int_a^bg\]

\seccion{6. Teorema fundamental y regla de Barrow}
\begin{formula}
\textbf{Teorema fundamental:} si \(f\) es continua,
\[A(x)=\int_a^xf(t)\,dt\ \Longrightarrow\ A'(x)=f(x)\]
\textbf{Regla de Barrow:} si \(F'=f\) en \([a,b]\),
\[\int_a^bf(x)\,dx=\bigl[F(x)\bigr]_a^b=F(b)-F(a)\]
\end{formula}
No hace falta \(+C\): se cancela al restar.
\begin{ejemplo}
\(A(x)=\int_0^x(t^2+1)\,dt\Rightarrow A'(x)=x^2+1\), sin calcular la integral.
\end{ejemplo}
\begin{ejemplo}
\[\int_1^3(3x^2-2x)\,dx=\bigl[x^3-x^2\bigr]_1^3=18-0=18\]
\[\int_0^1e^{2x}\,dx=\Bigl[\frac{e^{2x}}2\Bigr]_0^1=\frac{e^2-1}2\approx3{,}19\]
\end{ejemplo}

\begin{ejemplo}
Función a trozos: \(f(x)=x\) en \([0,1]\) y \(f(x)=1\) en \([1,2]\).
\[\int_0^2f(x)\,dx=\int_0^1x\,dx+\int_1^2 1\,dx=\frac12+1=\frac32\]
\end{ejemplo}

\seccion{7. Cálculo de áreas}
\begin{formula}
\[A=\int_a^b|f(x)|\,dx,\qquad A=\int_a^b|f(x)-g(x)|\,dx\]
\end{formula}
\begin{enumerate}
  \item Busca los \textbf{cortes} con el eje (\(f=0\)) o entre curvas (\(f=g\)).
  \item Divide \([a,b]\) en esos puntos y estudia el \textbf{signo} en cada trozo (o qué curva va por encima).
  \item Integra cada trozo y suma los \textbf{valores absolutos}.
\end{enumerate}
\begin{center}
\begin{tikzpicture}
\begin{axis}[mini,width=6.6cm,height=3.8cm,xmin=-0.4,xmax=2.6,ymin=-0.4,ymax=4.8,title={Área entre \(y=2x\) e \(y=x^2\)}]
\addplot[name path=f,domain=-0.2:2.4]{2*x};
\addplot[name path=g,domain=-0.2:2.2,dashed]{x^2};
\addplot[fill=acento!25,draw=none] fill between[of=f and g,soft clip={domain=0:2}];
\end{axis}
\end{tikzpicture}
\end{center}
\begin{ejemplo}
Cortes: \(x^2=2x\Rightarrow x=0,\ x=2\). En \(x=1\), \(2x=2>1=x^2\).
\[A=\int_0^2(2x-x^2)\,dx=\Bigl[x^2-\frac{x^3}3\Bigr]_0^2=\frac43\ \text{u}^2\]
\end{ejemplo}
\begin{ejemplo}
\(f(x)=x^2-4x+3\) en \([0,3]\): corta en 1 y 3; \(f>0\) en \((0,1)\), \(f<0\) en \((1,3)\). Con \(F=\frac{x^3}3-2x^2+3x\):
\[A=\Bigl|F(1)-F(0)\Bigr|+\Bigl|F(3)-F(1)\Bigr|=\frac43+\frac43=\frac83\ \text{u}^2\]
(\(\int_0^3f=0\): las partes se compensan; no es el área.)
\end{ejemplo}

\begin{ejemplo}
Recinto entre \(y=4-x^2\) y el eje \(X\): cortes en \(x=\pm2\) y \(f\ge0\) entre ellos.
\[A=\int_{-2}^{2}(4-x^2)\,dx=\Bigl[4x-\frac{x^3}3\Bigr]_{-2}^{2}=\frac{16}3+\frac{16}3=\frac{32}3\ \text{u}^2\]
\end{ejemplo}

\seccion{8. Aplicación en economía}
Si se conoce la tasa marginal, la integral da la \textbf{variación total}:
\[B(q_2)-B(q_1)=\int_{q_1}^{q_2}B'(q)\,dq\]
\begin{ejemplo}
\(B'(q)=40-2q\) miles de euros por centena. De 5 a 10 centenas: \(\int_5^{10}(40-2q)\,dq=\bigl[40q-q^2\bigr]_5^{10}=125\) miles de euros.
\end{ejemplo}

\seccion{9. Cómo elegir el método}
\begin{tabular}{@{}p{3.2cm}p{4.6cm}@{}}
\textbf{Si ves\ldots} & \textbf{Prueba\ldots}\\ \hline
Suma de potencias & Linealidad y fila de la potencia\\
Función interior y su derivada & Casi inmediata (ajusta números)\\
\(\frac{u'}u\), salvo constante & Logaritmo \(\ln|u|\)\\
Producto o potencia desarrollable & Desarrollar y simplificar\\
Cociente de polinomios & Dividir o separar términos\\
Raíz con interior y su derivada & Cambio de variable\\
\(x\cdot e^x\), \(\ln x\) & Por partes (ampliación)
\end{tabular}

\seccion{10. Problema encadenado}
\(f'(x)=2x-4\) y \(f(0)=3\). Halla \(f\) y el área entre su gráfica y el eje \(X\) en \([0,3]\).
\begin{enumerate}
  \item \(f(x)=\int(2x-4)\,dx=x^2-4x+C\); \(f(0)=C=3\Rightarrow f(x)=x^2-4x+3\).
  \item Es la función del ejemplo de áreas: \(A=\frac83\) u\(^2\).
\end{enumerate}

\begin{ojo}
\begin{itemize}
  \item Olvidar \(+C\) en las indefinidas (no en las definidas).
  \item \(\int\frac1x\,dx\) no es \(\frac{x^0}0\): es \(\ln|x|\).
  \item Dividir mal en el caso lineal: \(\int e^{3x}dx=\frac{e^{3x}}3\), no \(3e^{3x}\).
  \item Calcular un área con la integral de una función que cambia de signo.
  \item En el cambio de variable, dejar \(dx\) mezclado con \(u\).
\end{itemize}
\end{ojo}

\end{multicols}
\end{document}
